% Transcription — fonds Alexandre Grothendieck, shelfmark 64, batch 3 (pages 41-60).
% Established from the facsimile archives/batches/64/batch-03.pdf (a copy of 64.pdf
% fetched through the project relay, no cover sheet: PDF page k = archivists' page 20(N-1)+k),
% per the skill transcribe-grothendieck. The folder is 156 pages in eight batches.
%
% Pass: Opus 5.5 (claude-opus-5-5), 2026-09-27 — first pass, unchecked against the pages by a human.
%
% Pages skipped: 43 (a cover sheet in his hand, « (25 p.) II Structures octaédrales dans P1 »;
%   no \page{}, its words become the section heading). Page 42 carries a single line and is kept.
% His pagination: a new run starts on p. 44 = his p. 1; he numbers the even archive pages
%   (44 = 1, 46 = 2, ..., 60 = 9). Page 41 closes the preceding part (it ends « TSVP »).
% What the batch is about: p. 41-42, End(V) for V of rank 2 over F_2 (Aut^-(V) and Aut^+(V)
%   plus 0 as two commutative planes, trace form, sl(V)). From p. 44, part II « Structures
%   octaédrales dans P1 »: § 1 faithful actions of V ~ D_2 on a P^1-bundle X/S (char != 2),
%   axes Delta_1..3, normal form (4)-(5), the octahedral configuration Delta = {0, oo, ±1, ±i},
%   Gamma = Aut(X, Delta) ~ S_4 = Aut^+_oct(Delta), Theorem: bijections between sub-group
%   schemes ~ D_4, ~ S_4 and regular octahedral configurations; § 2 relative curves of type
%   (0,4): I ⊂ X étale of degree 4, V_I, Sigma_J the conic sum x_i^2 = 0, torsors K,
%   X* = K ∧_V Sigma*_{J,a}; § 3 relation with elliptic curves: E torsor under E_0 with an
%   involution sigma, I = E^sigma torsor under 2E_0, X = E/{1,sigma}, Sigma = E_0/{1,sigma_0},
%   Delta = 4-torsion points modulo ±, diagram (29).
% Notation fixed once for the batch: underlined sheaf/scheme objects \underline{V}, \underline{I},
%   \underline{J}, \underline{K}, \underline{M}; \mathbb{D}_2 (or \mathbb{D}_4, as written) for
%   the Klein group; \Sigma_{\underline{J}}, D_{\underline{J}}, \Sigma^*_{\underline{J}};
%   {}_2E_0 for 2-torsion; \underline{\operatorname{Aut}}, \underline{\operatorname{Pic}}.
%   His « icosaédrale » where octaédrale is meant (pp. 51, 52, 59) is kept as written.
% \ill{}: 96. \uncertain{}: 43.

\documentclass[11pt,a4paper]{article}
\input{../preamble/grothendieck}

\folder{64}
\batch{3}
\pages{41}{60}
\dating{[à partir de 1980- 1982]}
\watermark{Édition de démonstration}
\foldertitle{Modules courbes elliptiques en caractéristiques ≠2 : notes manuscrites (s.d.)}

\begin{document}

\page{41}
\textbf{Prop.} $V$ vectoriel de \uncertain{dim}~2 sur $\mathbb{F}_2$, donc
$V \simeq V_J(\mathbb{F}_2)$, où $J = V^*$ ens.\ à trois éléments,
$\operatorname{Aut}(V) \xrightarrow{\;\sim\;} \mathfrak{S}_J$. Posons
$\operatorname{Aut}^{\pm}(V) \simeq \mathfrak{S}_J^{\pm}$.
\uncertain{Considérons} $\operatorname{Aut}(V) \subset \operatorname{End}(V) = M$
\struck{\ill{}} (alg.\ d'Azumaya).
\note{les exposants de $\operatorname{Aut}(V)$ et de $\mathfrak{S}_J$ dans
« Posons \ldots » sont de petits traits superposés, lus $\pm$ d'après la suite
($\operatorname{Aut}^-$, $\operatorname{Aut}^+$). Sous « $\operatorname{Aut}(V)$ »
de la ligne précédente, une accolade renvoie au mot « Considérons », lui-même
surchargé. La lettre $M$ après $\operatorname{End}(V) =$ est suivie de lettres
biffées ; c'est la notation reprise p.~42.}

Alors :

1) $\operatorname{Aut}^-(V) = V'^{*}$ \struck{\ill{}}, i.e.
\[
  V' \subset \operatorname{End}(V), \qquad
  V' = \operatorname{Aut}^-(V) \cup \{0\}
\]
est un sous-espace vectoriel de rg~2 commutatif, canoniquement isomorphe à $V$,
en associant à $x \in V^*$ la transposition $\tau_x \in \operatorname{Aut}^-(V)$.
\note{le $V'$ de « $V'^{*}$ » est écrit par-dessus un $W$ ; l'égalité
$V' = \operatorname{Aut}^-(V) \cup \{0\}$ est écrite sous $V'$, reliée par
un signe $=$ vertical. Devant « $x \in V^*$ », un signe $\forall$ biffé.}

2) $\operatorname{Aut}^+(V) = W^*$, i.e.
\[
  W \subset \operatorname{End}(V), \qquad
  W = \operatorname{Aut}^+(V) \cup \{0\}
\]
est un sous-espace vectoriel de rang~2 commutatif,

3) $V'$ et $W$ sont orthogonaux l'un à l'autre pour la forme trace, qui sur
$V'$ se réduit à la forme bil.\ sym.\ unique qui est alternée et non dégénérée
\[
  \Bigl(\mathrm{Tr}\, xy =
  \begin{cases}
    0 & \text{si } x = y, \text{ ou } x \text{ ou } y = 0\\
    1 & \text{sinon i.e.\ } x, y \in V'^{*},\ x \neq y
  \end{cases}\Bigr),
\]
et sur $W$ est non alternée ($\mathrm{Tr}\,\rho^2 = \mathrm{Tr}\,\rho^{-2} = 1$,
$\mathrm{Tr}\,1^2 = 0$) et bien sûr non dégénérée.
\note{dans la marge gauche, en regard de ces deux formules, les matrices
$\begin{pmatrix} 0 & 1 \\ 1 & 0 \end{pmatrix}$ et
$\begin{pmatrix} 1 & 0 \\ 0 & 1 \end{pmatrix}$ : les matrices de Gram de la
forme trace sur $V'$ et sur $W$.}

4) $V'$ \struck{\ill{}} \ill{} \uncertain{que} $\mathbb{F}_2(\Gamma)$-module est
\uncertain{isom.}\ à $V$ et irréductible, tandis que $W$ est une extension non
scindée
\[
  0 \to \underset{\lambda \mapsto \lambda \cdot 1}{\mathbb{F}_2} \longrightarrow W
  \xrightarrow{\;\mathrm{Tr}\;} \mathbb{F}_2 \to 0 .
\]
\marginal{\emph{NB} $\mathbb{F}_2 \cdot 1$ est son propre orthogonal}
\note{la lettre lue $\Gamma$ dans « $\mathbb{F}_2(\Gamma)$ » est une grande
potence et reste douteuse. En regard du 4), dans la marge gauche, un
griffonnage au crayon, non lu. En bas à droite : « TSVP ».}

\page{42}
5) On a
\[
  \mathfrak{sl}(V) \overset{\text{déf}}{=}
  \operatorname{Ker}\bigl(M \xrightarrow{\;\mathrm{Tr}\;} \mathbb{F}_2\bigr)
  = V' + \mathbb{F}_2 \cdot 1_V .
\]
\note{le reste de la page est blanc. Le symbole lu $\mathfrak{sl}(V)$ est
écrit comme « $\int\!\ell(V)$ ».}

\section{II. Structures octaédrales dans $\mathbb{P}^1$}
\note{titre porté seul, de la main de l'auteur, sur la page~43, qui ne
porte rien d'autre que, en tête à gauche, l'indication « (25~p.) ».}

\subsection{1. Opération de groupes $\simeq \mathbb{D}_2$ sur la droite projective, et configuration octaédrale.}

\page{44}\note{p.~1 de l'auteur, chiffrée en tête de page.}
Soit $V$ un groupe qui est un espace vectoriel de rang~2 sur $\mathbb{F}_2$,
et considérons un \struck{e droite pro\ill{}} fibré en droites projectives $X$
sur un schéma de base $S$. On s'intéresse aux opérations, fidèles sur chaque
fibre, de $V$ sur $X$, i.e.\ aux monomorphismes
\[
  (1) \qquad V_S \hookrightarrow \underline{\operatorname{Aut}}_S(X) = G
\]
Choisissons une base $e_1, e_2$ de $V$ sur $\mathbb{F}_2$, alors la donnée d'un
homom.
\[
  V_S \longrightarrow \underline{\operatorname{Aut}}_S(X) \quad \text{i.e.}\quad
  V \longrightarrow \operatorname{Aut}_S(X)
\]
qui soit mono sur $\mathbb{F}_2 e_1$ et sur $\mathbb{F}_2 e_2$, \uncertain{équivaut à}
celle de \add{\uncertain{deux} automorphismes d'ordre deux de $X$, ou encore, de}
« deux » « axes » (i.e.\ diviseurs étales de degré~2) $\Delta_1$, $\Delta_2$ sur
$X$, la condition de commutation étant que $\Delta_2$ est stable par $\sigma_1$,
ou \uncertain{encore} $\Delta_1$ stable par $\sigma_2$.
\struck{La condition \ill{} $\sigma_1$ donné i.e.\ $\Delta_1$ donné, \ill{}
\ill{} que \ill{} ; \ill{} $\Delta_2$ \ill{} la donnée \ill{} mono signifie
$\sigma_2$ \ill{} que la donnée \ill{} $X/\{1, \sigma_1\} \to$ d'une section de}
\[
  (2) \qquad \Delta_1 \cap \Delta_2 = \emptyset \quad \text{i.e.}\quad
  \Delta_2 \subset X \setminus \Delta_1
\]
et stable par $\sigma_1$ (\uncertain{i.e.} $\Delta_1$ donné), la donnée de
\struck{\ill{}} $\sigma_2$ i.e.\ de $\Delta_2$ équivaut à : celle d'une section
de \struck{$X \setminus \Delta_1$ \ill{}} $(X \setminus \Delta_1)/\{1, \sigma_1\}$
sur $S$ ($\Delta_2$ étant alors l'image inverse de cette section dans
$X \setminus \Delta_1$). Considérons le troisième élément
\[
  (3) \qquad e_3 = e_1 + e_2
\]
il lui correspond un axe $\Delta_3$, et $\Delta_1$, $\Delta_2$, $\Delta_3$ sont
mutuellement disjoints.
\marginal{on est en car.\ $\neq 2$ ; $\Delta_1 = X^{\sigma_1}$,
$\Delta_2 = X^{\sigma_2}$}
\note{le passage biffé, sur quatre lignes, n'est lu qu'en partie. La note
marginale, écrite de biais dans la marge gauche en regard de « deux axes »,
définit les axes comme lieux des points fixes.}

\page{45}
La donnée de $\Delta_1$ est unique \ill{} localement près
\struck{\ill{}} --- si on prend une trivialisation de $\Delta_1$,
\struck{on peut prendre} la situation $(X, \Delta_1)$ i.e.\ \struck{\ill{}} est
isom.\ localement à $(\mathbb{P}^1_S, \{0, \infty\})$, $\sigma_1$ étant donné
par $z \mapsto -z$, donc $X \setminus \Delta_1$ étant
$\mathbb{P}^1_S \setminus \{0, \infty\} = \mathbb{G}_{m\,S}$, et le quotient
$(X \setminus \Delta_1)/\{1, \sigma_1\}$ s'identifiera à $\mathbb{G}_{m\,S}$,
l'homom.\ de passage au quotient étant


\begin{tikzcd}
X \setminus \Delta_1 \arrow[r] \arrow[d, "\wr"'] & (X \setminus \Delta_1)/\{1, \sigma_1\} \arrow[d, "\wr"] \\
\mathbb{G}_{m\,S} \arrow[r, "\lambda \mapsto \lambda^2"] & \mathbb{G}_{m\,S}
\end{tikzcd}

La donnée d'une section de $(X \setminus \Delta_1)/\{1, \sigma_1\}$ équivaut donc
à : celle d'une section $\lambda$ de $\mathbb{G}_{m\,S}$ --- et loc.\ pour la
top.\ étale, on peut poser $\lambda = \xi^2$, \struck{et quitte à \ill{} : la
section \uncertain{remplacer}} i.e.\ à la donnée d'une section $\xi$ de
$X \setminus \Delta_1$ sur $S$. On peut (quitte à remplacer l'iso.\
$X \simeq \mathbb{P}^1_S$ par un autre, savoir en multipliant par $1/\xi$)
supposer que $\xi$ est la section $1$, de sorte que la situation standard est
\[
  (4) \qquad
  \begin{cases}
    \Delta_1 = \{0, \infty\}, & \Delta_2 = \{1, -1\} \\
    \sigma_1(z) = -z & \sigma_2(z) = \dfrac{1}{z}
  \end{cases}
\]
(NB $\sigma_2$ est l'unique homographie qui échange $0$ et $\infty$, et fixe
$1$ et $-1$). On trouve donc
\note{la première ligne de la page n'est lue qu'en partie : le mot entre
« unique » et « localement près » est illisible, de même que le mot biffé
qui suit.}

\page{46}\note{p.~2 de l'auteur.}
\[
  (5) \qquad \sigma_3(z) = -\frac{1}{z} \qquad
  \Delta_3 = \mathbb{V}(z^2 + 1)
\]
($= \{i, -i\}$ si $\exists$ racine primitive \ill{} sur $S$)
On trouve donc

\textbf{Proposition.} Loc.\ pour la top.\ étale, une représentation
monomorphique de $\mathbb{D}_4 = \mathbb{F}_2 \times \mathbb{F}_2$ dans
$\operatorname{Aut}(\mathbb{P}^1_S)$ (sur une base $S$ à car.\ résid.\ $\neq 2$)
est donnée par (4) (où $\sigma_1$, $\sigma_2$ sont les \uncertain{automorphismes}
correspondant aux éléments de la base canonique $(e_1, e_2)$ de $\mathbb{D}_4$
sur $\mathbb{F}_2$) --- qui donne bien une représentation monomorphique.
\note{l'indice de $\mathbb{D}_4$ est bien un 4 dans les deux occurrences de
cette page, alors que le titre du §\,1 (p.~44) écrit
$\mathbb{D}_2$ ; nous transcrivons tel quel.}

Posons
\[
  (6) \qquad \Delta = \Delta_1 \amalg \Delta_2 \amalg \Delta_3 \qquad
  \bigl(= \{0, \infty, 1, -1, i, -i\} \text{ dans le cas normalisé}\bigr)
\]
on appelle $\Delta$ la \emph{configuration octaédrale} associée à
l'opération de $V$ sur \struck{\ill{}} $X$ --- par descente \uncertain{comme}
diviseur étale plat, elle est définie (comme diviseur étale de degré~6 de
$X/S$) dès qu'on a une opération monomorphique \add{(sur chaque fibre)} d'un
groupe $V \simeq \mathbb{D}_{2}$ sur $X$, ou plus généralement
d'un schéma en groupes $\underline{V}$ sur $S$ loc.\ isom.\ au schéma en
groupes constant. \struck{\ill{}} De façon générale, on appelle
\emph{configuration octaédrale} dans $X/S$ un diviseur étale de degré~6
de $X/S$, qui soit
\note{l'indice du $\mathbb{D}$ dans « $V \simeq \mathbb{D}_{2}$ », ajouté
au-dessus de la ligne, est peu net.}

\page{47}
loc.\ isomorphe à $\{0, \infty, 1, -1, i, -i\}$ dans $\mathbb{P}^1_S$, ou ce qui
revient au même, qui soit loc.\ isom.\ à la configuration octaédrale associée à
une op.\ \add{monom.}\ d'un $\underline{V}$ sur $X$. Nous allons voir maintenant
que cette donnée est essentiellement unique.

Pour ceci, pour un diviseur octaédral \struck{\ill{}}
\[
  \Delta \subset X
\]
considérons le schéma en groupes
\[
  (7) \qquad \Gamma = \underline{\operatorname{Aut}}_S(X, \Delta),
\]
il est immédiat (du fait que \uncertain{$\Gamma$}
$\subset \underline{\mathfrak{S}}_\Delta = \underline{\operatorname{Aut}}_S(\Delta)$,
\add{$\Delta$ ayant ses fibres} de degré $\geq 3$) que $\Gamma$ est
\emph{net} sur $S$, \uncertain{je dis} que $\Gamma$ est étale sur $S$, et
loc.\ isom.\ au schéma en groupes $(\mathfrak{S}_4)_S$. Plus précisément, je vais
définir sur $\Delta$ \add{(comme schéma des sommets)} une structure
d'\uncertain{octaèdre} \struck{\ill{}} \uncertain{de} schéma en octaèdres
combinatoires, de telle façon qu'on ait
\[
  (8) \qquad \Gamma = \underline{\operatorname{Aut}}^{+}_{\mathrm{oct.}}(\Delta)
  \qquad \text{(en tant que sous-schémas de } \underline{\operatorname{Aut}}_S(\Delta)\text{)}
\]
Pour ceci, on \add{note qu'une} \struck{\ill{}} structure octaédrale
sur un ens.\ à six éléments $\Delta$ est bien déterminée quand on connaît le
groupe de ses automorphismes \add{directs} (en tant que sous-groupe de
$\mathfrak{S}_\Delta$), car
\note{les deux insertions au-dessus de « une structure d'oct\ldots » et de
« on \ldots structure octaédrale » ne sont lues qu'en partie ; le début de la
ligne « \ldots\ schéma en octaèdres combinatoires » est peu sûr.}

\page{48}\note{p.~3 de l'auteur.}
les « diagonales » $\Delta_1$, $\Delta_2$, $\Delta_3$ de $\Delta$ sont
déterminées \add{(en tant que parties à deux éléments de $\Delta$)}, comme les
\struck{par la condition que leurs stabilisateurs soit} classes pour la
relation d'équivalence $R$ définie par l'application $s \mapsto \Gamma_s$
($s$ et $s'$ sont dans une même classe ssi ils ont le même groupe d'isotropie,
qui est \add{cyclique} d'ordre~4), et la structure polyédrale est déterminée en
termes de cette décomposition, \uncertain{les} arêtes étant les
\add{12}\ parties à deux éléments distinctes des $\Delta_i$, et les \add{8}\
faces étant les parties à trois éléments \emph{sections} de
$\Delta \to \Delta/R$. Donc la structure octaédrale sur $\Delta$ revient à la
structure \struck{de} définie par $R$, relation d'équivalence de degré~2, ou
encore par celle du groupe $\Gamma$ ($\simeq \mathfrak{S}_4$) opérant sur
$\Delta$.

Cette parenthèse combinatoire montre que la structure octaédrale canonique sur
$\Delta \subset X$ est nécessairement unique (puisqu'on connaît son schéma en
groupes d'autom.), ce qui montre que la question d'existence est locale
(\uncertain{p.t.}), donc OPS $X = \mathbb{P}^1_S$,
$\Delta = \{0, \infty, 1, -1, i, -i\}$. Le \struck{stabilisateur}
\add{sous-groupe} de \struck{$0, \infty$} \struck{$\Gamma$} qui fixe $0$ et
$\infty$ est isomorphe visiblement à : $\mu_4$ \struck{\ill{}}
$\simeq (\mathbb{Z}/4\mathbb{Z})_S$ (générateur $i$), \struck{et par} et de même
\ill{} (\ill{} d'homographies) pour $(1, -1)$, $(i, -i)$. Par contre, la
\note{« OPS » : on peut supposer. Le mot « \ill{} » après « et de même »
commence par « générat\ldots » et se lit peut-être « générateurs » ; celui
qui précède « d'homographies » n'est pas lu.}

\page{49}
groupe des autom.\ de $(X, \Delta)$ qui fixe $\infty$ et $-1$ donc
\add{(des autom.)} de la forme $z \mapsto az + b$ avec $-a + b = -1$ i.e.
\struck{$z \mapsto a(z+1) - 1 = az + (a-1)$}
\[
  z \longmapsto az + (a-1)
\]
est plus petit, en fait si l'autom.\ envisagé n'est pas l'identité, il doit
induire une permutation sans pts fixes de $(0, 1, i, -i)$, qui est donc d'ordre
2 ou quatre, donc \struck{on} $a^4 = 1$ \struck{ou $a^4 = 1$}. \add{mais faisant
$z = 0$ on doit avoir} \struck{i.e.\ $a \in \{0, 1, i, \ldots$}
\[
  a - 1 \in \{1, i, -i\} \quad \text{i.e.}\quad a \in (2, 1+i, 1-i)
\]
\marginal{donc $a^4 = \pm 4$ \uncertain{qui n'est pas égal à $1$ si car.\ $\neq 3$ et $5$}}
faisant $z = 1$ on trouve
\[
  2a - 1 \in \{0, i, -i\} \quad \text{i.e.}\quad
  a \in \Bigl\{\tfrac12, \tfrac12 (i+1), \tfrac12 (1-i)\Bigr\}
\]
\marginal{donc $a^4 = \frac14$ ou \uncertain{faisant} $a^4 = -1$
\uncertain{qui n'est pas égal à $1$} (car on est en car.\ $\neq 2$)}
faisant $z = i$, on trouve
\[
  a\underbrace{(i+1)}_{1 : (\frac{1-i}{2})} - 1 \in \{0, 1, -i\}
  \quad \text{i.e.}\quad a \in \Bigl\{\tfrac{1-i}{2}, 1-i, -i\Bigr\}
\]
et de même, faisant $z = -i$, on trouve $a \in \bigl\{\frac{1+i}{2}, 1+i, i\bigr\}$.
Donc l'intersection des stabilisateurs de $\infty$, $-1$ est $\{1\}$, et par
raison d'homogénéité il en est de même pour toute partie à deux élts de
$\Delta$, distincte des $\Delta_1$, $\Delta_2$, $\Delta_3$. Donc on voit que la
structure octaédrale $\Delta = \Delta_1 \amalg \Delta_2 \amalg \Delta_3$ est
\uncertain{connue} par $\Gamma$, et uniquement déterminée par le diviseur
$\Delta$. Donc $\Gamma \hookrightarrow \underline{\operatorname{Aut}}_{\mathrm{oct}}(\Delta)$,
\marginal{et même $\Gamma \subset \underline{\operatorname{Aut}}^{+}_{\mathrm{oct}}(\Delta)$}
Pour montrer qu'\ill{} isomorphisme, il suffit de noter que le groupe des
autom.\ \add{directs} de la structure octaédrale est engendré par les trois
sous-groupes cycliques d'ordre~4 \struck{\ill{}} qu'il contient,
correspondants \uncertain{aux diagonales} : $\Delta_1$, $\Delta_2$, $\Delta_3$
--- or ceux-ci sont contenus dans $\Gamma$.
\note{les deux notes marginales sur $a^4$, écrites de biais dans la marge
gauche en regard des cas $z = 1$ et $z = i$, ne sont lues qu'en partie ;
chacune se rattache au cas qu'elle borde. Dans le cas $z = i$, le premier
élément de l'ensemble final est surchargé ; les trois valeurs se déduisent de
$a(i+1) \in \{1, 2, 1-i\}$. La note « et même $\Gamma \subset \ldots$ »,
marquée d'un trait vertical, renvoie à « Donc
$\Gamma \hookrightarrow \underline{\operatorname{Aut}}_{\mathrm{oct}}(\Delta)$ ».}

\page{50}\note{p.~4 de l'auteur.}
Soit $\Delta$ un diviseur étale de degré~6 dans $X/S$, muni d'une structure
octaédrale i.e.\ d'une \struck{\ill{}} \add{\ill{}} \add{opération involutive}
\struck{de degré~2}
\[
  \underline{a}_\Delta : \Delta \to \Delta
\]
\struck{$\Delta \to \Delta'$}
\struck{i.e.\ $\Delta'$ est fini étale de degré~3 sur $S$.} \add{correspondant}
: une opération libre de $\{1, \underline{a}_\Delta\}$ sur $\Delta/S$. On dit
que cette structure octaédrale est \emph{régulière} \struck{\ill{}} dans
$X$, si le groupe $\Gamma = \underline{\operatorname{Aut}}^{+}_{\mathrm{oct}}(\Delta)$
opère sur $X$, de façon compatible avec les op.\ sur $\Delta$ \add{(cette op.\
de $\Gamma$ sur $X$ est alors unique)}. \struck{\ill{}} Si \add{le diviseur}
$\Delta$ \add{est} fixé, on voit qu'il y a au plus une structure octaédrale sur
$\Delta$ qui soit régulière, et il y en a une ssi $\Delta$ est une
« configuration octaédrale ». Si $\Delta$ est une telle configuration, on vient
de voir qu'elle est \add{dirigée,} munie d'une \struck{unique} structure
octaédrale, régulière (et tout autom.\ de $X$ qui \struck{fixe} \add{respecte}
$\Delta$ \struck{fixe} \add{respecte} cette structure octaédrale).
Inversement, si $\Delta$ est muni d'une str.\ oct.\ rég., alors le sous-groupe
$\underline{V}$ de $\Gamma = \underline{\operatorname{Aut}}_{\mathrm{oct}}(\Delta)$
invariant abélien non trivial est un groupe loc.\ isom.\ à
$\mathbb{D}_{2\,S}$, opérant fidèlement sur $X$, et on voit aussitôt que la
\marginal{Il faut un petit argument pour montrer que toute structure oct.\
$\Delta$ \ill{} induite par une op.\ de \ill{} est droite : OPS que
l'automorphisme fixe les points de $\{0, \infty\} = \Delta_1$, et on trouve dans
les autom.\ de « \uncertain{rotation} » \ill{} par $\mu_4$, qui sont bien
directs\ldots}
\marginal{Si on a une configuration octaédrale régulière \struck{$\Delta$} dans
un fibré \add{en droites} $X$, alors le revêtement des orientations de
\struck{$\Delta$} $\Delta$ est canoniquement isomorphe au revêtement
$\mu_4^{*} = \operatorname{Spec} \mathcal{O}_S[i]/(i^2+1)$ de $S$\ldots}
\note{les deux notes marginales sont écrites en diagonale dans la marge
gauche, la première en haut de la page (en regard de la définition de la
structure octaédrale), la seconde plus bas, d'une encre plus noire ; leur
lecture est incomplète. Près de « Soit », en exposant, « $(X, \Delta)$ ».
Une insertion au-dessus de « cette op.\ de $\Gamma$ sur $X$ » porte encore
« \uncertain{définie par des projections} », d'une lecture peu sûre.
L'indice de $\mathbb{D}_{2\,S}$ est surchargé (un 4 corrigé en 2, ou
l'inverse). Le mot souligné et biffé en fin de la ligne « \ldots\ alors
unique). » n'est pas lu.}

\page{51}
configuration \struck{icosaédrale} \add{octaédrale} associée est bien celle de
départ\ldots

\emph{Théorème} Soit $X$ un fibré en droites projectives sur un schéma $S$,
de car.\ \struck{\ill{}} résiduelles $\neq 2$. Considérons le schéma en groupes
$G = \underline{\operatorname{Aut}}_S(X)$ (\emph{NB} la donnée de $X/S$
équivaut à celle de la « forme » $G$ de $\mathrm{PGL}(2)_S$ sur $S$), et les
ensembles suivants :

a) L'ens.\ $\mathcal{A}(X)$ des sous-schémas en groupes $V$ de $G$, loc.\ isom.\
à $(\mathbb{D}_4)_S$

b) L'ens.\ $\mathcal{B}(X)$ des sous-schémas en groupes \struck{\ill{}} $\Gamma$
de $G$, loc.\ isom.\ à $(\mathfrak{S}_4)_S$.

c) L'ens.\ des configurations icosaédrales régulières dans $X$, i.e.\ des
diviseurs étales $\Delta$ dans $X/S$, de degré~6, munis d'un antipodisme
$\underline{a}_\Delta : \Delta \to \Delta$ qui \struck{\ill{}} \uncertain{soit
sans pt fixe} sur chaque fibre géom., et tel que
$\underline{\operatorname{Aut}}^{+}_{\mathrm{oct}}(\Delta)$ opère sur $X$ de
façon compatible avec son opération sur $\Delta$.

On a alors un système transitif de bijections canoniques entre ces trois
ensembles, obtenues ainsi !
\note{en c), l'auteur écrit « icosaédrales » là où il vient de corriger
« icosaédrale » en « octaédrale » en tête de page ; nous transcrivons tel quel.
Dans $\underline{\operatorname{Aut}}^{+}_{\mathrm{oct}}(\Delta)$, le $\Delta$
est surchargé. Le passage biffé après « qui » (deux mots, dont
« \uncertain{définit} ») n'est pas lu.}

\page{52}\note{p.~5 de l'auteur.}
\begin{itemize}
  \item[a) $\to$ b)] $\Gamma = \underline{\operatorname{Norm}}_G \underline{V}$
  \item[b) $\to$ a)] $\underline{V}$ est l'unique sous-schéma en groupes
    \add{(distingué)} fini étale de $\Gamma$, \struck{de} de degré~4,
    \struck{distingué}
  \item[c) $\to$ b)] $\Gamma = \underline{\operatorname{Aut}}_S(X, \Delta)
    \xleftarrow{\;\sim\;} \underline{\operatorname{Aut}}^{+}_{S\text{-oct}}(\Delta)$
  \item[a) $\to$ c)] \struck{$\underline{V}^*$} Toute section $\sigma$ de
    $V^* = V \setminus (\text{section unité})$ définit une action \add{fidèle}
    de $(\mathbb{Z}/2\mathbb{Z})_S$ sur $X$, donc un « axe » $\Delta_\sigma$ de
    $X$, et $\Delta$ est \struck{la réunion} (loc.\ ét.) le \struck{\ill{}}
    diviseur somme de ses axes\ldots
\end{itemize}
\marginal{\emph{NB} $\Gamma/V \simeq \underline{\operatorname{Aut}}_{S\text{-gr}}(V)
\simeq \underline{\operatorname{Aut}}_S(V^*)$}

\emph{Corollaire 1} Pour un diviseur étale de degré~6 $\Delta$ dans $X$, il
existe au plus une structure icosaédrale régulière ; s'il en existe une, le
\struck{groupe des autom.} \add{schéma en groupes} de ses automorphismes
\add{directs} est isom.\ à $\underline{\operatorname{Aut}}_S(X, \Delta)$.

\emph{Corollaire 2} La catégorie des couples $(X, \Delta)$
\add{sur $S$} (i.e.\ $(X, \underline{V})$, $(X, \underline{\Gamma})$\ldots)
est équivalente : à celle des \add{octaèdres} \struck{combinatoires}
\add{$\widetilde{F}$} sur $S$, ou encore : à celle des revêtements finis étales
\struck{\ill{}} de degré~4 de $S$, en associant à $(X, \Delta)$ le revêt.\ des
« faces modulo antipodisme » de la structure octaédrale $\Delta$, ou ce qui
revient au même, celui des sous-$3$-groupes de Sylow de
$\underline{\mathfrak{S}}_\Delta$ \ldots, celui des \ldots\
$\underline{\operatorname{Aut}}^{+}_{\mathrm{oct}}(\Delta) \simeq
\underline{\operatorname{Aut}}_S(X, \Delta)$. La structure-type${}^{*}$
correspond à \struck{revêt.} $X \subset \mathbb{P}^2_{\mathbb{Z}[\frac12]}$,
$X = \mathbb{V}(x^2 + y^2 + z^2 = 0)$, avec l'op.\ de
$\{\pm\}^3/\mathrm{diag} \simeq \mathbb{V}_{\{0,1,\infty\}}(\mathbb{F}_2)$
donnée par les chgts de signe de $x$, $y$, $z$.${}^{*}$
\marginal{\add{munis d'un isom.} $\underline{\mathrm{Or}}_{\Delta/S} \simeq
\mu^{*}_{4S}$ (on dit qu'on s'est donné une \uncertain{$\mu_4^{*}$-orientation} sur
$S$)}
\marginal{\emph{NB} Ici $X$ est totalement anisotrope sur $\mathbb{R}$ i.e.\
n'a pas de pts réels sur $\mathbb{R}$ --- \ill{} le revêtement \ill{} du groupe
$\mathbb{V}_{\{0,1,\infty\}}(\mathbb{F}_2)$ de $\mathbb{P}^1$ \ill{} avec une
\uncertain{structure} \ill{}, \ill{}\ldots\
${}^{*}$ Le \uncertain{biface} distingué est formé des deux faces
$\{(0,1,i), (i,0,1), (1,i,0)\}$ et $\{(0,1,-i), (-i,0,1), (1,-i,0)\}$.}
\note{la page est très chargée. Les flèches a)~$\to$~b), etc., renvoient aux
trois ensembles du Théorème p.~51. La note « munis d'un isom.\
\ldots\ » est écrite de biais à droite, au-dessus de « couples $(X, \Delta)$ »,
auxquels un trait la rattache ; le symbole lu $\underline{\mathrm{Or}}$ est
surchargé. La note NB finale est écrite verticalement dans la marge gauche,
de bas en haut ; seules les formules et quelques mots s'y lisent. Dans
« celui des sous-$3$-groupes de Sylow de $\underline{\mathfrak{S}}_\Delta$ »,
le $\mathfrak{S}$ est surchargé ; les deux « \ldots » qui suivent recouvrent
chacun un groupe de mots non lu. Devant $\{\pm\}^3$ un symbole biffé
($\mathbb{F}$ ?).}

\subsection{2. Courbes relatives de type $(0,4)$.}

\page{53}
Soit $X$ un fibré en droites projectives sur $S$, et
\[
  (9) \qquad \underline{I} \subset X
\]
un diviseur relatif, étale de degré~4. Considérons
\[
  (10) \qquad \underline{\mathfrak{S}}_{\underline{I}} = \underline{\operatorname{Aut}}_S(\underline{I})
\]
qui est un schéma en groupes fini étale sur $S$, loc.\ isom.\ à
$(\mathfrak{S}_4)_S$. Il lui correspond un unique sous-groupe \add{distingué}
$\underline{V}_{\underline{I}}$, loc.\ isom.\ à $(\mathbb{D}_2)_S$,
\struck{de telle façon} Alors $\underline{I}$ est un torseur sous
$\underline{V}_{\underline{I}}$.

\emph{Proposition} L'opération de $\underline{V}_{\underline{I}}$ sur
$\underline{I}$ se prolonge (de façon évidemment unique) en une opération de
$\underline{V}_{\underline{I}}$ sur $X$.

\emph{NB} La donnée de $\underline{I}$, fini étale de degré~4 sur $S$,
équivaut à : celle de $\underline{V} = \underline{V}_{\underline{I}}$ sur $S$
(ou encore de $\underline{V}^{*}_{\underline{I}} = \underline{J}_{\underline{I}}$,
fini étale de degré~3 sur $S$), plus la donnée d'un torseur $\underline{I}$ sous
$\underline{V}$. Plonger ce torseur $\underline{I}$ dans un fibré en droites
projectives $X$ revient à : faire opérer fidèlement $\underline{V}$ sur $X$
(chose étudiée en §\,1), et de plus, de se
\note{l'indice de $(\mathbb{D}_2)_S$ est surchargé. Le symbole
$\underline{J}_{\underline{I}}$ est écrit comme un grand $\amalg$ ; la lecture
$J$ suit la notation $J = V^*$ du début du dossier.}

\page{54}\note{p.~6 de l'auteur.}
donner une section de $(X \setminus \Delta)/\underline{V}$, \ill{} prenant
\add{pour $\underline{I}$} l'image inverse de cette section dans
$X \setminus \Delta$.

Mais la donnée d'une structure octaédrale équivaut à : celle d'un
\struck{\ill{}} rev.\ étale de degré~4 de $S$ (à ne pas confondre avec le
$\underline{I}$ du début, \uncertain{mais bien} le $\underline{K}$), ou encore,
à celle de

a) un schéma en groupes $V$ sur $S$, loc.\ isom.\ à $(\mathbb{D}_2)_S$ --- ou
encore un revêtement fini étale $\underline{J}$ ($= \underline{V}^*$) de degré~3
de $S$, ou encore, à celle d'un \struck{\ill{}} fibré en droites projectives
$\Sigma = \Sigma_{\underline{J}}$ sur $S$, muni d'un diviseur \add{relatif}
étale $D_{\underline{J}}$ de degré~3 (\ill{} isom.\ à $\underline{J}$)

b) d'un torseur $\underline{K}$ sous $\underline{V}$, ou ce qui revient au même,
d'un revêtement \add{principal} $X^*$ de
$\Sigma^*_{\underline{J}} = \Sigma_{\underline{J}} \setminus D_{\underline{J}}$
de groupe $\underline{V}$ \struck{$H^0(S, R^1 f_*(V)$} \struck{\ill{} \ill{}
fibre géométrique} correspondant \ill{} :
\[
  H^0\bigl(S, \mathcal{H}^1(\Sigma^*_{\underline{J}}/S, \underline{V})\bigr)
  \simeq \operatorname{Hom}(\underline{V}_{\underline{J}}, \underline{V})
\]
\struck{\ill{}}
--- la correspondance entre $\underline{K}$ et $X$ étant obtenue par la
formule
\[
  (11) \qquad X^* = \underline{K} \wedge_V \Sigma^*_{\underline{J},a},
\]
(après « $X^* =$ », \struck{$\Sigma_{\underline{J},0} \ldots$})
où $\Sigma^*_{\underline{J},a}$ est le revêtement principal dit
\marginal{\emph{NB} le \uncertain{rev.} \uncertain{univ.}\ local des
$(\pi_1)$ \ill{} $(V_{\underline{J}})_{\Sigma^*_{\underline{J}}}$ \ill{}
$\Sigma^*_{\underline{J}}$ sur $S$ \ill{} \ill{} l'identité
$\underline{V}_{\underline{J}} \to \underline{V} = \underline{V}_{\underline{J}}$
canoniquement isom.}
\note{la note marginale est écrite en diagonale dans la marge gauche, en
regard de la formule en $H^0$ ; elle n'est lue qu'en partie. La ligne de
« fibre géométrique » est biffée d'un long trait. Dans le « \ill{} » de la
première ligne, un mot court (« ein » ?) ; avant « isom.\ à $\underline{J}$ »,
un mot non lu.}

\page{55}
$\Sigma^*_{\underline{J}}$ défini \struck{par} à partir \struck{des} \struck{fibrés}
\add{de la conique} \struck{\ill{}} d'équation homogène
\[
  (12) \qquad \sum_{i \in \underline{J}} x_i^2 = 0
\]
\add{sur} laquelle le groupe $V_{\underline{J}}$
\add{$\simeq \mathbb{F}_2^{\underline{J}}/\mathbb{F}_2$} opère \struck{\ill{}}
diagonalement par changement de signe sur les $x_i$, comme quotient de
$\Sigma_{\underline{J}}$, défini par l'équation homogène
\[
  \sum x_i = 0 .
\]
Donc la donnée de $(X, \underline{I})$ équivaut : celle de \struck{\ill{}}
$\underline{J}$ i.e.\ de $\Sigma_{\underline{J}}$, du torseur $\underline{K}$ sous
$V_{\underline{J}}$ i.e.\ du revêtement \struck{\ill{}} $X$ de
$\Sigma_{\underline{J}}$ \struck{\ill{}} de groupe $V_{\underline{J}}$
(« admissible » \struck{\ill{}} sur qu'il définit l'identité
$V_{\underline{J}} \to V_{\underline{J}}$ \add{\uncertain{\ill{} \ill{}}}),
\emph{et} d'une section \struck{\ill{}} de $\Sigma^*_{\underline{J}}$
\marginal{[l'hypothèse que $\underline{I}$ soit contenu dans un fermé
\ill{} \uncertain{fibres}]}
\ldots\ Quand $\underline{I}$ est \add{fixé et} de la forme $I_S$, avec
\add{\ldots} éléments fixés (i.e.\ fixé d'\ill{}), on a
$\underline{J} \simeq J_S$, et $\Sigma^*_{\underline{J}} = \Sigma_J$
\ldots\ aussi fixé. La donnée de $(X, I_S \hookrightarrow X)$ équivaut
donc : celle de la section $x$ de $\Sigma^*_J$, section de « pt \uncertain{uniforme} »,
le revêtement $X$ de $\Sigma^*_J$ de groupe
$(V_J)_S = (V)(J_S) = \underline{V}$ étant
\marginal{\uncertain{description} \ill{} loc.\ ?) : il \ill{} \ill{}
\ill{} dans $\check{\mathbb{P}}(p_*(\mathcal{O}_{\underline{J}}))$ la conique
d'équation $Q = 0$, où $Q(u \otimes) = \mathrm{Tr}\, u^2$
(\emph{NB} \uncertain{deux} $p : \underline{J} \to S$ \ill{}
\uncertain{morphisme} structural) --- description \ill{} de l'\ill{}, loc.\
ét.\ --- la description \ill{} \ill{} l'hyperplan (13) de
$\check{\mathbb{P}}(p_*(\mathcal{O}_{\underline{J}}))$, d'équation
$L(u) = \mathrm{Tr}\, u$, $L(\ldots) = 0$ \ldots}
\marginal{Relations entre $\underline{I}$, $\underline{K}$, $x$ :
(14) $\underline{I} \simeq \underline{K} \wedge_{\underline{V}}
(\Sigma^*_{\underline{J},a}/x)$. Ainsi, pour $\underline{J}$ et
$x \in \Gamma(\Sigma^*_{\underline{J}}/S)$ fixés, $\underline{I}$ et
$\underline{K}$ (ou encore, $\underline{I}$ et $X$) \ill{} \ill{} \ill{}
\ill{} de degré~4 \ill{} de $S$, \ill{} $\underline{K}$ \ldots}
\note{la page est chargée ; le corps, surtout vers le bas (« Quand
$\underline{I}$ est fixé \ldots »), n'est lu qu'en partie. Les deux notes
marginales sont écrites en diagonale dans la marge gauche, la première en
regard de (12)--(13), la seconde (qui porte le numéro (14)) en bas ; elles
sont d'une lecture très incomplète. La note entre crochets à droite est
reliée par une flèche à la section de $\Sigma^*_{\underline{J}}$. Le numéro
(13) n'apparaît que dans la marge.}

\page{56}\note{p.~7 de l'auteur.}
\uncertain{déterminé} par la structure supplémentaire
\[
  (15) \qquad X/x \xrightarrow{\;\sim\;} I_S \qquad (V_S\text{-isomorphisme})
\]
qui, en termes du $\underline{V}_S$-torseur $\underline{K}$ qui a servi à définir
$X$, s'écrit
\struck{(16) $\underline{K} \simeq (I_S) \wedge \ldots (\Sigma^*_{J,a}/x)$}
\[
  (16) \qquad \underline{K} \simeq (I_S) \wedge_{\underline{V}} (\Sigma^*_{J,a}/x)
  \qquad \underline{V} = (V_J)_S .
\]
\note{le premier mot de la page est lu « déterminé » sous réserve ; il
commence par « d\ldots isant ». La première formule (16) est biffée et
réécrite ; son indice sous $\wedge$ est surchargé. Le reste de la page est
blanc.}

\subsection{3: Relations avec courbes elliptiques \add{: le morphisme $f_2^t : \Sigma_{\underline{J}} \to \Sigma_{\underline{J}}$}.}

\page{57}
Supposons toujours les car.\ rés.\ de $S$ $\neq 2$. Soient $E_0$ une courbe
elliptique sur $S$, $E$ un torseur sous $E_0$ (de sorte que
$E_0 \simeq \underline{\operatorname{Pic}}^0_{E/S}$). Supposons donné dans $E$
un \add{$S$-}\struck{\ill{}} \uncertain{automorphisme}
\[
  (17) \qquad \sigma : E \to E
\]
tel que
\[
  (18) \qquad \sigma_{E_0} = -\mathrm{id}_{E_0}
\]
ce qui implique que loc.\ (ét.), identifiant $E$ à $E_0$ via une section
origine de $E$, on aura
\[
  \sigma x = a - x \qquad a \text{ section convenable}
\]
ce qui est d'ailleurs intrinsèque si on interprète $a$ comme une section de
$E \wedge_{E_0} E \simeq \underline{\operatorname{Pic}}^2_{E/S}$. On aura donc
\[
  \sigma^2 = \mathrm{id}_E
\]
Soit \struck{\ill{}}
\[
  (19) \qquad
  \begin{cases}
    \underline{I} = E^\sigma \quad \text{sous-schéma des pts fixes de } \sigma \\
    \underline{I} \subset E \\
    \phantom{\underline{I}} = \{x \in E \mid 2x = a\}
  \end{cases}
\]
on voit que $\underline{I}$ est un torseur sous le schéma en groupes
\[
  (20) \qquad \underline{V} = {}_2E_0 = \operatorname{Ker}(2\,\mathrm{id}_{E_0} : E_0 \to E_0),
\]
qui est justement un groupe \add{sur $S$} loc.\ isom.\ à $\mathbb{D}_{2\,S}$.
La donnée de $(E, \sigma)$, où $E$ est un schéma
\note{dans (18), l'exposant de $\sigma$ (un $2$ corrigé ?) est surchargé ;
la lecture $\sigma_{E_0}$ (l'automorphisme induit sur $E_0$) est celle que
donne l'indice. Le mot lu « automorphisme » après « un $S$- » est écrit
court et en partie biffé. Dans (19), l'égalité de la troisième ligne se
rattache à $\underline{I}$ par un signe $=$ vertical.}

\page{58}\note{p.~8 de l'auteur.}
propre et lisse sur $S$ à fibres géom.\ connexes de dim~1, genre~1, et
$\sigma$ un autom.\ induisant $-\mathrm{id}$ sur
$\underline{\operatorname{Pic}}^0_{E/S}$ \add{équivaut à : celle} de
$(E, a)$, où $a$ section de $\underline{\operatorname{Pic}}^2_{E/S}$
($\simeq E \wedge_{E_0} E$), ou encore à celle de $(E_0, \underline{I})$, où
$E_0$ est une courbe elliptique relative sur $S$, et où $\underline{I}$ est un
torseur sous \struck{${}_2E_0 \overset{\text{déf}}{=}$}
$\underline{V} \overset{\text{déf}}{=} {}_2(E_0)$. On aura
\[
  (21) \qquad
  \begin{cases}
    E = \underline{I} \wedge_{\underline{V}} E_0 \\
    \sigma \text{ déduit de la symétrie } x \mapsto -x \text{ de } E_0
    \text{ via la formule précédente}
  \end{cases}
\]
On aura une opération de $\underline{V} = {}_2E_0 \subset E_0$ sur $E$, et un
isom.\ canonique
\[
  (22) \qquad
  \begin{array}{c}
    E/\underline{V} \xrightarrow{\;\sim\;} E_0 \\
    \wr\!\wr \\
    E \wedge_{E_0} E \simeq \underline{\operatorname{Pic}}^2_{E/S}
  \end{array}
\]
isom.\ qui transforme la symétrie $\sigma$ sur $E$, \add{en}
$\sigma_0 : x \mapsto -x$ sur $E_0$.
\marginal{\emph{NB} L'opération de $\underline{V}$ sur $E$ commute à
l'action de $\sigma$}

Considérons les schémas quotients
\[
  (23) \qquad X = E/\{1, \sigma\}, \qquad \Sigma = E_0/\{1, \sigma_0\}
\]
ce sont des fibrés en \struck{\ill{}} \add{droites} projectives,
\add{(canoniquement \uncertain{isomorphes})}, \struck{et sur les} $V$ opère
\add{sur} $X$ \struck{et sur} $\Sigma$ via son action sur $E$, et on aura
\[
  (24) \qquad X/V \xrightarrow{\;\sim\;} \Sigma ,
\]
d'ailleurs le diviseur de ramification pour
\[
  (25) \qquad E \to X \qquad (\text{rev.\ quadratique})
\]
\note{le « $\wr\!\wr$ » de (22) rend un signe d'égalité vertical entre
$E_0$ et $E \wedge_{E_0} E$ ; un « 21 » est écrit à côté, peut-être un
renvoi à (21). L'insertion lue « canoniquement isomorphes », sous « fibrés
en droites », est peu nette ; son rattachement à $X$ et $\Sigma$ est
incertain.}

\page{59}
est $\underline{I}$, et on a un \struck{\ill{}}\add{\emph{monomorphisme}}
\[
  (26) \qquad \underline{I} \hookrightarrow X
\]
\struck{stable sous l'action} compatible avec l'action de $\underline{V}$, dont on
déduit facilement que c'est un \emph{monomorphisme} (car on vérifie que
$\underline{I}$ s'envoie dans l'\ill{} $X^*$ de $X$ où l'action de
$\underline{V}$ est libre, cf.\ plus bas (27)). Ainsi on est dans la situation
du §\,2, qui implique que $\Sigma \simeq \Sigma_{\underline{J}}$, où
$\underline{J} = \underline{V}^*$, que $\underline{I}$ est l'image inverse d'une
section bien déterminée \struck{\ill{}} de
$\Sigma^*_{\underline{J}} = \Sigma_{\underline{J}} \setminus D_{\underline{J}}$,
l'image inverse $\Delta$ de $\underline{J}$ dans $X$ étant \struck{\ill{}} la
configuration icosaédrale dans $X$, correspondant à l'action donnée de
$\underline{V}$ sur le \struck{\ill{}} fibré en droites projectives $X$.

En termes de l'action de \struck{$V$} $\underline{V}$ sur $X$,
$\underline{\Delta}$ se récupère comme le diviseur de ramification pour cette
action --- ensemblistement, il correspond donc aux pts (à val.\ dans des corps
alg.\ clos) dont le stabilisateur dans $V$ est $\neq 0$, i.e.\ il correspond :
aux pts $x \in E$ tels qu'il existe $\xi \in V = {}_2E_0$, $\xi \neq 0$, avec
\[
  x + \xi \in \{x, \sigma x = a - x\} \quad \text{i.e.} \quad
  x + \xi = a - x \quad \text{i.e.} \quad 2x + \xi = a
  \quad \text{d'où} \quad 4x = 2a \quad \text{mais } 2x \neq a
\]
\struck{i.e.\ \ldots}
\note{l'auteur écrit « configuration icosaédrale » où la p.~46
parlait de configuration octaédrale. Le mot devant $X^*$ (« l'o\ldots ») n'est
pas lu. Dans « $V = {}_2E_0$ », l'indice est surchargé (un 8 ?). La fin biffée
de la dernière ligne porte un $E$ et un indice $4$.}

\page{60}\note{p.~9 de l'auteur.}
\[
  x \text{ pt de } \underline{I} \wedge_{{}_2E_0} ({}_4E_0) \setminus \underline{I}
\]
On trouve ainsi
\[
  (27) \qquad \Delta = \Bigl(\underbrace{\underbrace{\underline{I} \wedge_{\underline{M}_0} \underline{M}}_{\text{degré } 16}
  \setminus \underbrace{\underline{I}}_{\text{degré } 4}}_{\text{degré } 12}\Bigr)
  \Big/ \underbrace{\{1, \sigma\}}_{\text{degré } 2}
\]
où on a \uncertain{posé maintenant}
\[
  (28) \qquad
  \begin{cases}
    \underline{M} \simeq {}_4E_0 \quad \text{schéma en groupes loc.\ isom.\ à } (\mathbb{Z}/4\mathbb{Z})^2_S \\
    \underline{M}_0 \simeq \underline{M} \otimes_{\mathbb{Z}/4\mathbb{Z}} \mathbb{Z}/2\mathbb{Z}
    \simeq 2\underline{M} = {}_2E_0 = \underline{V}
  \end{cases}
\]
Ainsi \add{(dans le cas $E = E_0$, disons)} les pts d'ordre~4 de $E_0$ se
groupent par \add{6} paires d'éléments symétriques, qui forment la
configuration octaédrale $\Delta$ dans $X$.

On trouve ainsi le diagramme (\uncertain{presque cartésien})


\begin{tikzcd}[column sep=large]
\underline{I} \subset E \arrow[r, "\text{étale}"] \arrow[d] & E_0 \ (\simeq E/\underline{M}_0) \arrow[d] \\
\underline{I} \hookrightarrow X \arrow[r] & \Sigma_{\underline{J}} \\
\Delta \arrow[u, hook] \arrow[r, dashed] & D_{\underline{J}} \arrow[u, hook]
\end{tikzcd}

(29) --- qui est cartésien au-dessus de
$\Sigma^*_{\underline{J}} = \Sigma_{\underline{J}} \setminus D_{\underline{J}}$.
\note{le diagramme (29) est redessiné : sur la page, les flèches verticales
$E \to X$ et $E_0 \to \Sigma_{\underline{J}}$ portent respectivement
« ramifié au-dessus des pts de $\underline{I}$ » et « ramifié au-dessus des
pts de $D_{\underline{J}} \cup t(S)$ », la flèche horizontale
$X \to \Sigma_{\underline{J}}$ « ramifié au-dessus des pts de
$D_{\underline{J}}$ » ; les deux inclusions de la dernière ligne sont notées
par un signe $\cup$ sous $X$ et sous $\Sigma_{\underline{J}}$. Dans (28),
l'indice de $\otimes$ et le $\mathbb{Z}/2\mathbb{Z}$ (écrit « $\mathbb{Z}_{/2}$ »)
sont peu nets ; le signe entre $2\underline{M}$ et ${}_2E_0$ est surchargé.}

\end{document}
